% =========================================================
\section{Capteurs, détecteurs et fins de course}

\subsection{Introduction}

Dans un système automatisé, la partie commande doit connaître à tout moment l'état de la partie opérative et de son environnement.

Pour cela, elle utilise des \textbf{capteurs} et des \textbf{détecteurs}, qui traduisent une information physique en un signal exploitable par l'API.

Parmi ces organes, les \textbf{fins de course} occupent une place importante : ils indiquent la position extrême d'un organe mobile.

\subsection{Définition d'un capteur}

Un capteur est un composant qui prélève une information sur une grandeur physique (position, présence, niveau, température, pression, etc.) et la transforme en un signal électrique exploitable par la partie commande.

\[
\text{Grandeur physique} \rightarrow \text{Capteur} \rightarrow \text{Signal électrique} \rightarrow \text{API}
\]

\subsection{Classification des capteurs}

On peut classer les capteurs selon plusieurs critères.

\subsubsection{Selon le type de signal}

\begin{itemize}
    \item \textbf{Capteurs TOR} (Tout Ou Rien) : délivrent un signal 0 ou 1 ;
    \item \textbf{Capteurs analogiques} : délivrent un signal proportionnel à la grandeur mesurée (0--10 V, 4--20 mA) ;
    \item \textbf{Capteurs numériques} : délivrent un code (binaire, Gray, BCD).
\end{itemize}

\subsubsection{selon le mode de détection}

\begin{itemize}
    \item \textbf{Avec contact} : fin de course mécanique ;
    \item \textbf{Sans contact} : détecteurs de proximité, photoélectriques, ultrasoniques.
\end{itemize}

\subsubsection{Selon la nature de la grandeur}

\begin{itemize}
    \item Détection de présence ou de position ;
    \item Mesure de niveau ;
    \item Mesure de température ;
    \item Mesure de pression ;
    \item Mesure de débit ;
    \item Mesure de vitesse ou de position angulaire.
\end{itemize}

\subsection{Les fins de course}

\subsubsection{Définition}

Un fin de course est un capteur mécanique qui détecte la position extrême d'un organe mobile (vérin, porte, tiroir, etc.).

Il est généralement équipé d'un contact \textbf{NO} ou \textbf{NC} et d'un levier ou d'un galet actionné par l'organe en mouvement.

\subsubsection{Principe de fonctionnement}

Lorsque l'organe mobile arrive en fin de course, il appuie sur le levier du fin de course, ce qui ouvre ou ferme le contact.

\[
\text{Organe mobile en butée} \rightarrow \text{Action sur le levier} \rightarrow \text{Contact ouvert ou fermé}
\]

\subsubsection{Types de fins de course}

\begin{itemize}
    \item Fin de course \textbf{à levier} ;
    \item Fin de course \textbf{à galet} ;
    \item Fin de course \textbf{à tige} ;
    \item Fin de course \textbf{à billes}.
\end{itemize}

\subsubsection{Utilisation}

Les fins de course servent à :

\begin{itemize}
    \item détecter la position d'un vérin (sorti/rentré) ;
    \item limiter la course d'un organe mobile ;
    \item arrêter un moteur en fin de course ;
    \item confirmer une position avant d'autoriser l'étape suivante d'un cycle.
\end{itemize}

\paragraph{Exemple :} dans un vérin pneumatique, deux fins de course indiquent respectivement la position \textbf{rentrée} et la position \textbf{sortie}.

\subsection{Les détecteurs de proximité}

\subsubsection{Définition}

Un détecteur de proximité est un capteur \textbf{sans contact} qui détecte la présence d'un objet à faible distance.

\subsubsection{Détecteur inductif}

Il détecte uniquement les objets \textbf{métalliques}.

\paragraph{Principe :} il crée un champ magnétique alternatif et détecte la perturbation provoquée par un objet métallique qui s'approche.

\paragraph{Avantages :}
\begin{itemize}
    \item pas d'usure mécanique ;
    \item cadence élevée ;
    \item insensible à la poussière et à l'humidité.
\end{itemize}

\paragraph{Utilisation :} détection de cames, de pignons, de pièces métalliques.

\subsubsection{Détecteur capacitif}

Il détecte les objets \textbf{conducteurs ou isolants} (plastique, verre, bois, liquides).

\paragraph{Principe :} il détecte la variation de capacité provoquée par l'approche d'un objet.

\paragraph{Utilisation :} détection de niveau de liquide à travers une paroi, détection de pièces plastiques.

\subsubsection{Détecteur photoélectrique}

Il utilise un \textbf{faisceau lumineux} (infrarouge ou visible) pour détecter un objet.

Trois modes existent :

\begin{itemize}
    \item \textbf{Barrage} : émetteur et récepteur séparés ; l'objet coupe le faisceau.
    \item \textbf{Réflex} : émetteur et récepteur dans le même boîtier, avec un réflecteur ; l'objet coupe le faisceau.
    \item \textbf{Proximité} : l'objet lui-même renvoie le faisceau vers le récepteur.
\end{itemize}

\paragraph{Utilisation :} détection de pièces, comptage, sécurité (barrière immatérielle).

\subsubsection{Détecteur à ultrasons}

Il émet des ondes ultrasonores et mesure le temps de retour de l'écho.

\paragraph{Utilisation :} détection d'objets transparents, mesure de distance, détection de niveau.

\subsection{Les capteurs de grandeur physique}

\subsubsection{Capteurs de niveau}

Ils détectent ou mesurent le niveau d'un liquide ou d'un solide dans un réservoir.

\begin{itemize}
    \item \textbf{Niveau TOR} : flotteur, électrode, détecteur capacitif ;
    \item \textbf{Niveau analogique} : capteur de pression, ultrasons, radar.
\end{itemize}

\subsubsection{Capteurs de température}

\begin{itemize}
    \item \textbf{Thermocouple} : mesure de hautes températures ;
    \item \textbf{Sonde PT100} : mesure précise et linéaire ;
    \item \textbf{Thermistance} : mesure économique.
\end{itemize}

\subsubsection{Capteurs de pression}

Ils mesurent la pression d'un fluide (air, eau, huile).

\subsubsection{Codeurs}

Un codeur est un capteur qui mesure une position angulaire ou une vitesse de rotation.

\begin{itemize}
    \item \textbf{Codeur incrémental} : délivre des impulsions ; il faut un comptage ;
    \item \textbf{Codeur absolu} : délivre directement la position sous forme de code.
\end{itemize}

\subsection{Caractéristiques d'un capteur}

Pour choisir un capteur, on tient compte de :

\begin{itemize}
    \item la \textbf{grandeur} à détecter ;
    \item le \textbf{type de signal} de sortie (TOR, analogique, numérique) ;
    \item la \textbf{distance de détection} ;
    \item la \textbf{cadence} de fonctionnement ;
    \item la \textbf{tension d'alimentation} (24 V DC, 230 V AC) ;
    \item le \textbf{type de sortie} (PNP, NPN, contact sec) ;
    \item l'\textbf{environnement} (poussière, humidité, température, vibrations) ;
    \item l'\textbf{indice de protection} (IP).
\end{itemize}

\subsection{Sorties PNP et NPN}

Les détecteurs de proximité électroniques possèdent souvent une sortie à transistor.

\begin{itemize}
    \item \textbf{Sortie PNP} : le transistor relie la sortie au \textbf{+24 V} lorsqu'il est actif ;
    \item \textbf{Sortie NPN} : le transistor relie la sortie au \textbf{0 V} lorsqu'il est actif.
\end{itemize}

\paragraph{Attention :} il faut respecter le type de sortie attendu par l'entrée de l'API. Un capteur PNP ne se raccorde pas comme un capteur NPN.

\subsection{Raccordement à l'API}

Les capteurs et détecteurs sont raccordés aux entrées de l'API.

\paragraph{Exemple de tableau des entrées :}

\begin{table}[h]
\centering
\begin{tabularx}{\textwidth}{|c|X|c|c|}
\hline
\textbf{Adresse} & \textbf{Organe} & \textbf{Type} & \textbf{Sortie} \\
\hline
I0.0 & Fin de course vérin rentré & TOR & Contact NC \\
\hline
I0.1 & Fin de course vérin sorti & TOR & Contact NO \\
\hline
I0.2 & Détecteur inductif pièce & TOR & PNP \\
\hline
I0.3 & Détecteur photoélectrique & TOR & PNP \\
\hline
I0.4 & Capteur de niveau bas & TOR & Contact NO \\
\hline
IW64 & Capteur de pression & Analogique & 4--20 mA \\
\hline
\end{tabularx}
\caption{Exemple de raccordement de capteurs}
\end{table}

\subsection{Exemple d'application : détection de position d'un vérin}

Un vérin pneumatique doit être commandé en va-et-vient.

Deux capteurs de position sont installés :

\begin{itemize}
    \item un capteur \textbf{vérin rentré} (fin de course) ;
    \item un capteur \textbf{vérin sorti} (fin de course).
\end{itemize}

Le programme de l'API utilise ces informations pour :

\begin{itemize}
    \item autoriser la sortie du vérin uniquement lorsqu'il est rentré ;
    \item autoriser la rentrée uniquement lorsqu'il est sorti.
\end{itemize}

\[
\text{Vérin rentré} \rightarrow \text{API} \rightarrow \text{Commande sortie}
\]

\[
\text{Vérin sorti} \rightarrow \text{API} \rightarrow \text{Commande rentrée}
\]

\subsection{Exemple d'application : comptage de pièces}

Un détecteur photoélectrique détecte le passage de pièces sur un convoyeur.

Chaque détection est comptée par l'API. Lorsque le compteur atteint une valeur déterminée, l'API arrête le convoyeur ou déclenche une alarme.

\[
\text{Impulsion détecteur} \rightarrow \text{Compteur API} \rightarrow \text{Action}
\]

\subsection{Bonnes pratiques}

\begin{itemize}
    \item Vérifier la compatibilité PNP/NPN avec l'API ;
    \item Respecter la tension d'alimentation du capteur ;
    \item Éviter les environnements trop perturbés (vibrations, poussière) ;
    \item Repérer clairement chaque capteur sur le schéma ;
    \item Utiliser un contact NC pour les sécurités ;
    \item Tester chaque capteur avant raccordement ;
    \item Protéger les câbles contre les perturbations électromagnétiques.
\end{itemize}

\subsection{Erreurs fréquentes à éviter}

\begin{itemize}
    \item Confondre capteur PNP et NPN ;
    \item Utiliser un capteur inductif pour détecter du plastique ;
    \item Dépasser la distance de détection nominale ;
    \item Négliger l'indice de protection dans un environnement humide ;
    \item Mal régler la sensibilité d'un détecteur photoélectrique ;
    \item Ne pas tenir compte du temps de réponse du capteur.
\end{itemize}

\subsection{Conclusion}

Les capteurs, détecteurs et fins de course sont les organes qui informent la partie commande sur l'état du système.

Leur choix dépend de la grandeur à détecter, du type de signal, de l'environnement et du type de sortie attendu par l'API.

Un bon raccordement et un repérage rigoureux sont essentiels pour la fiabilité du système automatisé.

\subsubsection*{À retenir}

\begin{quote}
Un capteur transforme une grandeur physique en signal exploitable par l'API. Un fin de course détecte la position extrême d'un organe mobile. Un détecteur de proximité fonctionne sans contact.
\end{quote}

\paragraph{Schéma fondamental :}

\[
\text{Grandeur physique} \rightarrow \text{Capteur} \rightarrow \text{Signal} \rightarrow \text{Entrée API}
\]